\chapter{总结与展望}

\section{全文总结}
% TODO: 总结全文工作
% 本文面向多用户多关键词对称可搜索加密场景，研究了在泄露受限隐私保护模型下引入轻量可验证机制的问题。
% 主要工作与贡献包括：
%
% 第一，提出了基于 Merkle Hash Tree 的可验证资格检验机制。
% 该机制以 MHT 替换传统 Bloom Filter 组织授权信息，
% 为资格检验结果生成可验证的成员性证明路径，
% 将资格检验从概率判定提升为可证明判断，
% 消除了假阳性导致的错误拒绝风险，
% 并使授权用户能够检测云服务器的选择性作弊行为。
%
% 第二，提出了基于嵌入式承诺的搜索结果正确性验证机制。
% 该机制在数据更新阶段利用承诺机制将校验信息嵌入既有 SSE 索引结构，
% 无需引入独立验证索引，验证开销成为原有索引操作的边际增量。
% 承诺的绑定性保证云服务器无法在搜索阶段伪造与正确结果不一致的校验信息，
% 验证过程在用户端完成，不依赖云服务器的诚实性。
%
% 实验评估表明，两个可验证机制引入后的额外开销处于可接受范围，
% 验证了方案的可行性与工程可落地性。

\section{未来工作展望}
% TODO: 讨论未来研究方向
% - 前向隐私与后向隐私在可验证框架下的联合实现
% - 更丰富的查询语义支持（析取查询、范围查询等）
% - 基于区块链的去中心化可验证搜索
% - 更强对抗模型下的安全性增强（如抗串谋）
% - 面向大规模数据的性能优化与工程实现
